% Based on: Poke-RRT paper.

\chapter{Poking: Dynamic Primitives for Sampling-Based Planning}
\label{ch:poking}

% TODO: Introduce poking as the entry point into dynamics-constrained
% manipulation.

\section{Motivating Scenario}
\label{sec:poking:motivation}

% TODO: Two arms with non-overlapping workspaces --- a setting that is
% unsolvable by pushing or pick-and-place alone.

\section{Poking as a Complementary Primitive}
\label{sec:poking:primitive}

\subsection{Two-Stage Mechanics}
\label{sec:poking:mechanics}

% TODO: Formalize the impulsive strike followed by free sliding under
% Coulomb friction.

\subsection{Position Relative to Quasi-Static Manipulation}
\label{sec:poking:context}

% TODO: Situate poking within the broader space of quasi-static
% manipulation paradigms.

\section{The Poke-RRT Algorithm}
\label{sec:poking:algorithm}

\subsection{Heuristic Strike-Point Selection}
\label{sec:poking:strikepoint}

% TODO: Describe the heuristic used to select strike points on the object.

\subsection{End-Effector Action Computation}
\label{sec:poking:action}

% TODO: Joint-space conversion and feasibility checks for the chosen
% strike action.

\subsection{Forward Propagation in a Kinodynamic RRT Tree}
\label{sec:poking:rrt}

% TODO: Use simulation as forward propagation within the kinodynamic RRT.

\section{Results}
\label{sec:poking:results}

% TODO: Approximately 3x faster than pushing, ~2x fewer actions, and
% higher success rates than push and grasp baselines.

\section{Reflection on the Four Design Axes}
\label{sec:poking:reflection}

% TODO: Evaluate Poke-RRT against the modeling burden, dimensionality,
% integration, and usability axes from Chapter~\ref{ch:framework}.
